\documentclass[11pt]{article}
\usepackage[margin=1in]{geometry}
\usepackage[T1]{fontenc}
\usepackage{lmodern,amsmath,amssymb,amsthm,booktabs,array,etoolbox}
% Build with XeLaTeX; the Japanese quotation uses the IPA Mincho font.
\usepackage{fontspec}
\newfontfamily\japanesefont{IPAMincho}
\usepackage[hyphens]{url}
\usepackage[colorlinks=true,linkcolor=blue,citecolor=blue,urlcolor=blue]{hyperref}
\newtheorem{theorem}{Theorem}[section]
\newtheorem{proposition}[theorem]{Proposition}
\newtheorem{lemma}[theorem]{Lemma}
\newtheorem{corollary}[theorem]{Corollary}
\theoremstyle{definition}
\newtheorem{definition}[theorem]{Definition}
\newcommand{\C}{\mathbb C}
\newcommand{\Q}{\mathbb Q}
\newcommand{\Z}{\mathbb Z}
\newcommand{\F}{\mathbb F}
\newcommand{\Hilb}{\operatorname{Hilb}}
\newcommand{\Frac}{\operatorname{Frac}}
\newcommand{\rank}{\operatorname{rank}}
\newcommand{\Sym}{\operatorname{Sym}}
\newcommand{\ini}{\operatorname{in}}
\newcommand{\LT}{\operatorname{LT}}
\newcommand{\wt}{\operatorname{wt}}
\newcommand{\SLP}{\mathrm{SLP}}
\newcommand{\degrees}{\{6,12,18,24,30,42\}}
\pagestyle{plain}
\setlength{\emergencystretch}{2em}
\raggedbottom
\title{The strong Lefschetz property for complex reflection coinvariants\thanks{The proofs, the certificate checkers and the text were produced by AI models (from OpenAI and Anthropic), working under the author's direction.
This work has not yet been reviewed by experts.}}
\author{P. Graham\\ \small\url{https://orcid.org/0009-0007-8493-1505}}
\date{October 26, 2026}
\hypersetup{pdfauthor={P. Graham},pdftitle={The strong Lefschetz property for complex reflection coinvariants}}
\begin{document}
\maketitle
\begin{abstract}
The ordinary coinvariant algebra over $\C$ of every finite
complex reflection group has the strong Lefschetz property: one linear
element has maximal-rank multiplication by every power in every degree.
We give a computer-assisted proof for $G_{34}$, the exceptional type
excluded by Wachi's earlier results; the other types follow from the work
of Harima--Watanabe, Harima--Migliore--Nagel--Watanabe, McDaniel, and Wachi.
The question is recorded in Wachi's 2013 conjecture on complex reflection
coinvariants. The proof uses the Conway--Sloane orbit invariants in the
coordinates developed by Oura--Sekiguchi. Exact finite-field polynomial
certificates, specialization, and a Hilbert-function comparison establish
the strong Lefschetz property for a relative coinvariant algebra. The
free-extension theorem then transfers it to the ordinary coinvariant
algebra, using Wachi's result for $G_{33}$.
\end{abstract}


\section{Introduction}
Let $V$ be a nonzero finite-dimensional complex vector space, and let
$W\subseteq\operatorname{GL}(V)$ be a finite group generated by complex
reflections. A complex reflection is a nonidentity linear transformation fixing a
hyperplane pointwise. Give
$R=\Sym(V^*)$ its ordinary polynomial grading, and let $R^W_+$ denote the
positive-degree invariants. The \emph{coinvariant algebra} is
\[
                         E_W=R/(R^W_+R).
\]
For example, the cyclic reflection group of order $m$ in one variable
has coinvariant algebra $\C[x]/(x^m)$. Multiplication by each power of
$x$ has the largest possible rank in every degree.

\begin{definition}
A finite-dimensional graded algebra $D=\bigoplus_{i\geq0}D_i$ has the
\emph{strong Lefschetz property} if there is an element $\ell\in D_1$
such that
\[
 \rank(\ell^k:D_i\longrightarrow D_{i+k})
                 =\min\{\dim D_i,\dim D_{i+k}\}
                       \qquad(i,k\geq0).
\]
Such an element is a strong Lefschetz element. Graded algebras below
need not be generated by their degree-one pieces.
\end{definition}

\begin{theorem}\label{thm:main}
For every irreducible finite complex reflection group $W$ on a nonzero
complex vector space, the ordinary coinvariant algebra $E_W$ has the
strong Lefschetz property.
\end{theorem}

The product reduction in Corollary~\ref{cor:all} extends this statement
to reducible groups and representations with fixed directions.

The contribution here is the $G_{34}$ case. Wachi proved the other
exceptional nonreal cases of rank at least three, including $G_{33}$
\cite{Wachi25,WachiCode}; Wachi's author statement explicitly excludes
$G_{34}$. We use that $G_{33}$ result as a literature premise. The
precise source qualification is stated in Section~\ref{sec:trust}.
McDaniel and Smith jointly announced the Bott--Samelson embedding work
in 2016 \cite[pp.~5, 10--11]{BIRS16}. Smith's talk asserted an SLP
deduction for groups generated by order-two reflections, a class that
includes $G_{34}$ \cite[p.~6]{BIRS16}. The same report records
McDaniel's warning about inheritance \cite[McDaniel's talk, p.~5]{BIRS16}.
McDaniel--Smith \cite[Proposition~5.9]{McDSmith} exhibit a type-$A_3$
Bott--Samelson map without the inheritance property, although both rings
have SLP. Thus the embedding alone does not supply the required
multiplication ranks. This work supplies those ranks for $G_{34}$ by
relative-coinvariant certificates and free extension, rather than by
inheritance along a Bott--Samelson embedding. The cited noninheritance
example does not disprove the announced SLP conclusion for $G_{34}$.
The report also records a McDaniel--Smith--Wachi breakout lemma giving
sufficient conditions for SLP of the base in a Leray--Hirsch decomposition,
and Smith's use of it for $G_{25}$ and $G_{26}$
\cite[\S4.4, p.~10]{BIRS16}. This is an earlier base-SLP strategy close
to the free-extension step here; Wachi's 2013 Lemma~36 already treats
$G_{24}$ through $G_{27}$ \cite{WachiLecture}.
The difficulty is to obtain all multiplication ranks, rather than only
a nonzero highest power of a linear form. Instead of treating the full
$G_{34}$ coinvariant algebra directly, we prove the required ranks on
a smaller relative algebra through an auxiliary invariant quotient.
The free-extension method is due to Harima--Watanabe and
McDaniel--Chen--Iarrobino--Macias Marques
\cite{HW07,MCCM}. The orbit invariants are due to Conway--Sloane
\cite{CS}, and their useful coordinate descriptions and subgroup model
are supplied by Oura--Sekiguchi \cite{OS}.

\subsection*{Earlier cases and the question}
Wachi states Conjecture~37 in the proceedings of Wachi's November~29, 2013
lecture \cite[Conjecture~37, proceedings p.~155, PDF p.~16]{WachiLecture},
\cite{WachiProgram}. The original Japanese sentence is:
\begin{quote}
\japanesefont 複素鏡映群の余不変式環は、強レフシェッツ性を持つ。
\end{quote}
\noindent\emph{Translation:} ``The coinvariant ring of a complex
reflection group has the strong Lefschetz property.''
The sentence does not restrict the group to be irreducible.
Theorem~\ref{thm:main} treats the irreducible factors;
Corollary~\ref{cor:all} supplies the full assertion.
This is the earliest dated public presentation identified here; the
original online upload date of the notes is not established. Those notes
already list the five remaining cases $G_{29},G_{31},G_{32},G_{33},G_{34}$.
Wachi's
March~2014 abstract \cite[Functional Analysis, talk~12, p.~78]{WachiAbstract},
in the March~3 collection, also records the question with five
exceptional types remaining. The question is thus dated to its public
2013 presentation, rather than to the 2025 article.
The Shephard--Todd classification \cite{ST} reduces
Theorem~\ref{thm:main} to the cases in Table~\ref{tab:coverage}.
Wachi's lecture, Theorem~4 and Tables~1--3, is also a convenient source
for this classification and the degrees used below.

\begin{table}[ht]
\centering\small
\caption{Earlier results used for types other than $G_{34}$.}
\label{tab:coverage}
\begin{tabular}{@{}p{.25\textwidth}p{.69\textwidth}@{}}
\toprule
Types & Result and exact location \\
\midrule
Cyclic groups & The one-variable example above, including the trivial
rank-one quotient. \\[3pt]
Type $A$ and real reflection groups & McDaniel \cite[Theorem~1]{McD};
the proof in \S5 includes $H_3$ and $H_4$. Extension of scalars from
$\mathbb R$ to $\C$ preserves multiplication ranks. \\[3pt]
$G(r,p,n)$, $r\geq2$, $p\mid r$ & Harima--Watanabe
\cite[Notation and Remark~6.3, Examples~6.4 and~6.8]{HW06}, with $s=r/p$;
the invariant presentation is explained below. \\[3pt]
$G_4,\ldots,G_{22}$ & These have rank two. Harima--Migliore--Nagel--Watanabe
\cite[Proposition~4.4]{HMNW} prove SLP for homogeneous Artinian
quotients in two variables in characteristic zero. \\[3pt]
$G_{23},G_{28},G_{30},$\newline$G_{35},G_{36},G_{37}$ & These are
$H_3,F_4,H_4,E_6,E_7,E_8$, respectively, and are included in
McDaniel's real-group theorem. \\[3pt]
$G_{24},\ldots,G_{27},$\newline$G_{29},G_{31},G_{32},G_{33}$ &
Wachi \cite{Wachi25,WachiCode}. Lecture Lemma~36 already treats the
first four; the author repository lists all eight proved types and
leaves $G_{34}$. \\
\bottomrule
\end{tabular}
\end{table}

For clarity, the monomial-family citation concerns the ordinary
coinvariant algebra in the standard grading. Put $s=r/p$ and let $e_i$
be the $i$th elementary symmetric polynomial. Its presentation is
\begin{equation}\label{eq:monomial}
 \frac{\C[x_1,\ldots,x_n]}
 {(e_1(x_1^r,\ldots,x_n^r),\ldots,
    e_{n-1}(x_1^r,\ldots,x_n^r),(x_1\cdots x_n)^s)}.
\end{equation}
To see this, the diagonal ratios in $G(r,p,n)$ force invariant monomial
exponents to coincide modulo $r$, and the allowed single-coordinate
diagonal elements force that common residue to be a multiple of $s$.
After extracting a power of $t=(x_1\cdots x_n)^s$, permutation
invariance gives a symmetric polynomial in $x_i^r$. Its last elementary
symmetric generator is $t^p$. Thus the invariant ring is generated by
$e_1,\ldots,e_{n-1},t$; these are algebraically independent since
the polynomial ring is integral over its invariant subring and has
dimension $n$. This proves \eqref{eq:monomial}. For $p=1$, Example~6.4
of \cite{HW06} applies with $s=r$; for $p>1$, Example~6.8 applies with
$0<s<r$. Rank one is immediate. The case $r=1$ is type $A$, with its
fixed summand killed by its linear invariant.

For Weyl groups, McDaniel's introduction uses Borel's identification
with flag-variety cohomology and the hard Lefschetz theorem
\cite[Introduction]{McD}. Stanley's earlier work applies hard Lefschetz
to Weyl-group order structures \cite{Stanley80}. McDaniel also credits
Maeno--Numata--Wachi and Numata--Wachi for the noncrystallographic cases;
see \cite[Introduction and \S5]{McD}. Yoshida characterizes Lefschetz
elements for complex reflection groups containing a binary polyhedral
subgroup of index two \cite{Yoshida}; this related rank-two work is not
needed for the existence theorem supplied by \cite{HMNW}.
Related work of McDaniel--Smith
\cite{McDSmith} studies embeddings into Bott--Samelson algebras.
Such an embedding alone does not transfer SLP: their Proposition~5.9
exhibits the missing inheritance condition even in type $A_3$.
The argument here instead supplies SLP of both the base and fiber of
a free extension. It asserts neither an inheritance theorem for
arbitrary subalgebras nor an implication from maximal height to SLP.

\subsection*{Proof overview}
Use the subgroup chain
\[
                   H=G(3,3,4)\ \subset\ K=G_{33}\ \subset\ W=G_{34}.
\]
The subgroup $K$ fixes the linear form $z=(x_5-x_6)/2$. Let $A$ be the
relative coinvariant algebra for $K\subset W$. The ordinary $W$
coinvariant algebra is a free extension of $A$ with the ordinary $K$
coinvariant algebra as fiber. The latter has SLP by Wachi. It remains
to prove that $z$ is a strong Lefschetz element of $A$.

The invariant ring of $H$ is an explicit polynomial ring. Its quotient
by the positive $W$-invariants, denoted $C$, is a graded free
$A$-module with a known fiber $B$. For every $k$, this makes the
Hilbert function of $C/(z^k)$ at least the function predicted by
maximal-rank multiplication on $A$, multiplied by the Hilbert series of $B$.
Exact polynomial identities over $\F_{1009}$ give a monomial-ideal
upper bound. A finite monomial enumeration shows that the upper and
lower bounds agree for every required $k$. Specialization goes in
the needed direction because nonzero minors modulo $1009$ remain
nonzero in characteristic zero.

Section~\ref{sec:model} defines the model and its integral coordinates.
Section~\ref{sec:finite} states the finite facts certified by the
programs. Sections~\ref{sec:specialize} and~\ref{sec:free} prove the
bridges from those facts to SLP over $\C$. The final section describes
the certificate format, reproduction, and the trusted boundary.

\section{The group model and integral invariants}\label{sec:model}
This section fixes the characteristic-zero objects to which the modular
inputs must belong. Let $\omega=e^{2\pi i/3}$ and
$\theta=1+2\omega=\omega-\omega^2$. The following vectors in
$\Z[\omega]^6$ form a set $\mathcal V$:
\begin{align*}
 &\pm\theta(\omega^a e_i-\omega^b e_j),
       &&1\leq i<j\leq6,\quad a,b\in\{0,1,2\},\\
 &\pm(\omega^{a_1},\ldots,\omega^{a_6}),
       &&a_i\in\{0,1,2\},\quad \sum_i a_i\equiv0\pmod3.
\end{align*}
Here $e_i$ is the $i$th coordinate vector. The two sets have $270$ and
$486$ elements, respectively. This is the Conway--Sloane construction
\cite[Theorem~10]{CS} in the form described in \cite[\S2]{OS}.
For a vector $n\in\mathcal V$, its Hermitian squared norm is $6$,
and the corresponding reflection is
\[
                   s_n(u)=u-\frac{\overline n^{\,T}u}{3}\,n.
\]
In the notation of Oura--Sekiguchi, use the normals
\[
\begin{array}{c|cccccc}
 &P_1&P_2&P_3&Q_1&R_1&R_2\\ \hline
n&\theta(e_2-e_3)&\theta(e_3-e_4)&\theta(e_4-e_5)&
\theta(e_1-e_2)&\theta(e_1-\omega^2e_2)&(1,1,1,1,1,1).
\end{array}
\]
Let $W$ be generated by all six reflections, let
$K=\langle P_1,P_2,Q_1,R_1,R_2\rangle$, and let
$H=\langle P_1,P_2,Q_1,R_1\rangle$.
The identification of $W$ with $G_{34}$ and $K$ with $G_{33}$ is
\cite[\S\S2,6.1]{OS}; $H$ acts as $G(3,3,4)$ on the first four
coordinates and fixes the last two. The invariant degrees are
\begin{equation}\label{eq:degrees}
\begin{array}{c|l}
 W&(6,12,18,24,30,42)\\
 K&(1,4,6,10,12,18)\\
 H&(3,4,6,9,1,1).
\end{array}
\end{equation}
These are literature premises, also listed in
\cite[Tables~2--3]{WachiLecture}. The degree-one invariant of $K$
accounts for its fixed line.

The exact group check verifies preservation of $\mathcal V$, the
reflection matrices, inclusions, and orders
\begin{equation}\label{eq:orders}
             |W|=39\,191\,040,\qquad |K|=51\,840,
                         \qquad |H|=648.
\end{equation}
Its permutation action on $\mathcal V$ is faithful: six vectors in
that set have a determinant with nonzero norm $8748$ and hence span.
The subgroup $K$ fixes $e_5-e_6$; unitarity implies that it fixes its
dual coordinate $z$. The checker verifies these assertions by exact
Eisenstein-integer arithmetic and a faithful finite permutation
representation. Agreement of orders by itself is not used to identify
the named groups or their invariant degrees.

Put $y=(x_5+x_6)/2$, $z=(x_5-x_6)/2$, and define
\[
 a=\sum_{i=1}^4x_i^3,\qquad b=x_1x_2x_3x_4,\qquad
 c=\sum_{i=1}^4x_i^6,\qquad d=\sum_{i=1}^4x_i^9.
\]
Their weights, in the order $a,b,c,d,y,z$, are $(3,4,6,9,1,1)$.
Let
\[
 D=\bigl(\Z[\omega,1/42!]\bigr)_{(1009,\omega-374)},
                   \qquad L=\Frac(D)=\Q(\omega).
\]
The integer $1009$ is prime, and
$374^2+374+1=139\cdot1009$. Thus $D$ has the residue map
$\omega\mapsto374$ to $\F_{1009}$. All denominators through $42!$
are units.

\begin{lemma}\label{lem:Hinv}
The invariant ring of $H$ over $D$ is
$D[x_1,\ldots,x_6]^H=D[a,b,c,d,y,z]$, a polynomial ring with the
weights above. The same assertion holds over $\F_{1009}$ and over $L$.
\end{lemma}
\begin{proof}
The diagonal cube-root ratios force the first four exponents of an
invariant monomial to coincide modulo three. The differences of
distinct cube roots are units in $D$, so this follows coefficient by
coefficient. Every such monomial has a unique expression as $b^r$
times a monomial in $x_1^3,\ldots,x_4^3$, with $0\leq r<3$.
The permutation subgroup fixes $b$ and permutes these four cubes.
The fundamental theorem of symmetric polynomials over $\Z$ therefore
gives the elementary generators $e_1,e_2,e_3,e_4$ in the cubes, with
$e_4=b^3$. Their first three values are
\[
        e_1=a,\qquad e_2=(a^2-c)/2,\qquad
                            e_3=(a^3-3ac+2d)/6.
\]
This proves generation by $a,b,c,d$. For independence, split a
purported polynomial relation by the exponent of $b$ modulo three.
The three summands have disjoint exponent residues in the original
variables. Each would give a relation in the independent elementary
symmetric functions $e_1,e_2,e_3,e_4$, so all vanish. Finally, the
change $x_5=y+z$, $x_6=y-z$ is invertible since $2$ is a unit.
The same argument works in the residue field and in $L$.
\end{proof}

For $q\in\degrees$, define the unscaled orbit sum
\begin{equation}\label{eq:orbit}
                 f_q(x)=\sum_{v\in\mathcal V}(v^Tx)^q.
\end{equation}
The group check proves invariance: $\mathcal V$ is stable under the
generators and under conjugation, hence also under their transposes
and inverses. Thus either equivalent dual-action convention preserves
the sums. By Lemma~\ref{lem:Hinv}, each $f_q$ has a unique expression
$F_q\in D[a,b,c,d,y,z]$. Denote its reduction modulo the maximal
ideal by $I_q$, and put
\[
           S=\F_{1009}[a,b,c,d,y,z],\qquad
                         I=(I_6,I_{12},I_{18},I_{24},I_{30},I_{42}).
\]

The committed modular inputs are checked against \eqref{eq:orbit}
by an exhaustive coefficient calculation. For the second vector
family, its character sum is $486$ when all six exponent residues
modulo three agree and zero otherwise. The checker tests all $3^6$
residue tuples and expands every possible surviving exponent tuple;
it sums the first family on each two-coordinate support. It then
substitutes $x_5=y+z$, $x_6=y-z$ and compares all coefficients with
the substituted $I_q$. Multinomial coefficients use factorials at most
$42$. Lemma~\ref{lem:Hinv} makes the modular substitution injective,
so this proves that the six inputs are the stated reductions, rather
than a presentation fitted by point evaluations.

\section{The finite certificate facts}\label{sec:finite}
This section gives the exact finite assertions used by the proof.
Write $H_D(t)=\sum_j\dim(D_j)t^j$ for the Hilbert series of a finite
graded algebra over a field. Introduce the following explicit integer
polynomials:
\begin{align}
 h(t)&=\frac{(1-t^{24})(1-t^{30})(1-t^{42})}
                   {(1-t)(1-t^4)(1-t^{10})}
             =\sum_{j=0}^{81}h_jt^j,\label{eq:h}\\
 b(t)&=\frac{(1-t^{10})(1-t^{12})(1-t^{18})}
                   {(1-t)(1-t^3)(1-t^9)},\label{eq:b}\\
 Q_k(t)&=\sum_{j\geq0}\max(0,h_j-h_{j-k})t^j,
                     \qquad h_j=0\text{ for }j\notin[0,81].\label{eq:Q}
\end{align}
Their meaning as characteristic-zero Hilbert series will be established
in Section~\ref{sec:free}.

Use weighted degree reverse lexicographic order on $S$, with variables
$a,b,c,d,y,z$ in that order and $z$ last. The certificate gives
homogeneous polynomials $G_1,\ldots,G_{534}$ and polynomials $T_{qj}$.
Write $\LT(g)$ for the leading monomial of a nonzero polynomial $g$,
and $\ini(U)$ for the monomial ideal generated by the leading monomials
of all nonzero elements of an ideal $U$.
Let $J=(\LT(G_1),\ldots,\LT(G_{534}))$.

\begin{proposition}[Certified finite facts]\label{prop:finite}
The six modular inputs defined above and the certificate satisfy:
\begin{enumerate}
\item Every $G_j$ is nonzero and
$G_j=\sum_{q\in\degrees}I_qT_{qj}$ coefficient by coefficient.
The declared leading exponents agree with those reconstructed from
the polynomials and the specified order.
\item The monomial ideal $J$ contains the pure powers
$a^2,b^3,c^3,d^4,y^{68},z^{82}$. Its complement consists of exactly
$60\,480$ monomials, all of weighted degree at most $108$.
\item For every $1\leq k\leq81$, the complete Hilbert series satisfies
\[
                          H_{S/(J,z^k)}(t)=Q_k(t)b(t).
\]
\end{enumerate}
\end{proposition}

These are finite exact assertions, verified by the checkers distributed
with this paper. Two arithmetic implementations verify part~(1).
Two separate monomial enumerations verify parts~(2)--(3), using the
same checked leading exponents but separate traversal methods.
There is a finite bound before any Hilbert-series comparison: each
standard exponent is below the corresponding pure-power exponent.
Every monomial outside that box lies in $J$. Thus enumeration proves
the absence of a tail as well as agreement of the displayed
coefficients; the degree $108$ bound is an output of the complete
cover, not an assumed stopping point.

For one enumeration, every standard monomial is uniquely
$u z^i$, where $z$ does not divide $u$ and $0\leq i<L_u$.
The quotient by $z^k$ has the basis with
$0\leq i<\min(k,L_u)$. A separate enumeration traverses the entire
bounded standard set directly. Both calculate \eqref{eq:Q} by integer
coefficient arithmetic. Further checks certify a Gr\"obner basis and
the converse inclusion of ideals, but only the identities and the
monomial upper bound in Proposition~\ref{prop:finite} are needed below.

\section{From the residue field to the invariant ideal}\label{sec:specialize}
We now justify the characteristic-zero upper bounds and prove that
the chosen orbit sums generate the whole $W$-invariant ideal. These
steps do not assume flatness or an integral presentation for $R^K$.

\begin{lemma}[Specialization]\label{lem:specialize}
Put $\widetilde C=L[a,b,c,d,y,z]/(F_q:q\in\degrees)$. For every
$j\geq0$ and $k\geq1$,
\[
 \dim_L(\widetilde C/(z^k))_j
       \leq\dim_{\F_{1009}}(S/(I,z^k))_j
       \leq\dim_{\F_{1009}}(S/(J,z^k))_j.
\]
The same inequalities hold without adjoining $z^k$.
\end{lemma}
\begin{proof}
In weighted degree $j$, list the finitely many monomials. Make a
matrix over $D$ whose columns are all products of the $F_q$ and of
$z^k$ with monomials of the complementary degree. Its cokernel is
the degree-$j$ quotient. Every minor nonzero in the residue field is
already nonzero over $L$. Hence the characteristic-zero rank is at
least the residue rank, which proves the first inequality. This does
not lift a modular zero relation.

By Proposition~\ref{prop:finite}(1), $G_j\in I$, so
$J\subseteq\ini(I)$. Also $z^k$ is itself a leading monomial in
$I+(z^k)$. Therefore
$J+(z^k)\subseteq\ini(I+(z^k))$. The standard monomials of the
initial ideal are a vector-space basis of the corresponding quotient,
which gives the second inequality. This basis fact follows directly
by leading-term reduction: nonstandard monomials reduce to smaller
ones, and a nontrivial relation among standard monomials would have
a leading term in the initial ideal. Positivity of the weights makes
the order well-founded. No equality of the initial ideals is needed.
\end{proof}

\begin{proposition}\label{prop:invariant-generation}
Over $\C$, the orbit sums generate the invariant ring:
\[
                   R^W=\C[f_6,f_{12},f_{18},f_{24},f_{30},f_{42}],
                \qquad R=\C[x_1,\ldots,x_6].
\]
\end{proposition}
\begin{proof}
Proposition~\ref{prop:finite}(2) makes $S/J$ finite-dimensional.
Lemma~\ref{lem:specialize} without $z^k$ makes $\widetilde C$
finite-dimensional. After extending scalars to $\C$ and using
Lemma~\ref{lem:Hinv}, this says that $R^H/(f_q)$ is
finite-dimensional. The ring $R$ is finite over $R^H$: each element
is integral by its finite orbit polynomial, and $R$ is generated by
finitely many integral elements. Consequently $R/(f_q)$ is finite
as well.

Thus the six $f_q$ form a homogeneous system of parameters of the
six-variable polynomial ring. A polynomial ring is Cohen--Macaulay,
and such a parameter system is a regular sequence; see
\cite[Tags~00ND and~02JN]{Stacks}. Apply the local statement at
the homogeneous maximal ideal; a nonzero graded kernel remains
nonzero there. Regularity gives
\[
             H_{R/(f_q)}(t)=\frac{\prod_{q\in\degrees}(1-t^q)}{(1-t)^6},
           \qquad \dim R/(f_q)=\prod_{q\in\degrees}q=39\,191\,040.
\]

Let $P=\C[f_q:q\in\degrees]$. Homogeneous lifts of a vector-space
basis of $R/(f_q)$ generate $R$ over $P$, by induction on degree:
subtract the corresponding lift, express the remainder using the
positive-degree $f_q$, and apply the induction to its coefficients.
Thus $R$ is finite over $P$. Equal dimensions show that the six
generators of $P$ are algebraically independent. Indeed, the kernel
of the map from a polynomial ring in six variables would otherwise
be a nonzero prime and lower the dimension.

The same lifts give a surjection from a graded free $P$-module onto
$R$. The displayed regular-sequence Hilbert series shows that the
source and target have equal Hilbert series, so its kernel is zero.
Its rank is $\prod q=|W|$, by \eqref{eq:orders}, and therefore
$[\Frac R:\Frac P]=|W|$.
Artin's fixed-field theorem
\cite[Tag~09I3]{Stacks} gives
$[\Frac R:(\Frac R)^W]=|W|$ because the action is faithful.
Moreover, $\Frac(R^W)=(\Frac R)^W$: multiply the denominator of
an invariant rational function by its other group translates to make
an invariant denominator, and then its numerator is invariant too.
Since $P\subseteq R^W$, the equal field degrees imply
$\Frac P=\Frac(R^W)$.
Every element of $R^W$ is integral over $P$, because $R$ is finite
over $P$, and lies in this common field. A polynomial ring is
integrally closed, so $R^W=P$.
\end{proof}

In particular, after scalar extension, $\widetilde C$ is precisely
$C=R^H/(R^W_+R^H)$. Extension from $L$ to $\C$ preserves ranks
and dimensions in each finite graded piece. The specialization
bound of Lemma~\ref{lem:specialize} therefore applies to this
characteristic-zero invariant quotient.

\section{The free module and the Lefschetz comparison}\label{sec:free}
We relate the auxiliary quotient to the relative algebra on which
$z$ will be Lefschetz. All rings in this section are over $\C$.
In addition to $C$, put
\[
\begin{array}{ll}
 A=R^K/(R^W_+R^K),& F=R/(R^K_+R),\\[2pt]
 B=R^H/(R^K_+R^H),& E=R/(R^W_+R).
\end{array}
\]
Here $F$ is the ordinary $G_{33}$ coinvariant algebra; its fixed
linear coordinate is killed in the quotient. The Hilbert series
from the degrees in \eqref{eq:degrees} give the following values.
The top degree means the largest degree with a nonzero component.

\begin{center}
\begin{tabular}{@{}lrr@{}}
\toprule
Algebra & Dimension & Top degree\\
\midrule
$A$ & $756$ & $81$\\
$F$ & $51\,840$ & $45$\\
$E$ & $39\,191\,040$ & $126$\\
$B$ & $80$ & $27$\\
$C$ & $60\,480$ & $108$\\
\bottomrule
\end{tabular}
\end{center}

For completeness, each displayed relative quotient is a complete
intersection in the appropriate polynomial invariant ring. It is
finite because the smaller-group invariant ring is finite over the
larger-group one. Its parameters therefore form a regular sequence,
by the same Cohen--Macaulay argument as above. Taking the ratio of
the products $1-t^q$ gives
\begin{equation}\label{eq:Hilbert}
                H_A(t)=h(t),\qquad H_B(t)=b(t),\qquad
                                H_C(t)=h(t)b(t).
\end{equation}
In particular $A$ and $F$ have symmetric Hilbert functions:
replacing $t$ by $t^{-1}$ in their complete-intersection formulas
reverses the coefficients about their respective top degrees.

\begin{lemma}\label{lem:free}
There is a graded $A$-module isomorphism $C\cong A\otimes_\C B$
that intertwines multiplication by $z$. Consequently, for every
$k\geq1$,
\[
                     H_{C/(z^k)}(t)=H_{A/(z^k)}(t)b(t).
\]
\end{lemma}
\begin{proof}
The reflection-coinvariant free-extension theorem
\cite[Proposition~3.1]{MCCM} gives $A\longrightarrow E$
with fiber $F$, since $R^K$ is polynomial. The required injection
also has a direct description: if a $K$-invariant belongs to the
ideal $R^W_+R$, average its coefficients in a $W$-invariant
generator expression over $K$. This proves
$(R^W_+R)\cap R^K=R^W_+R^K$.

Choose a graded linear section $s:F\longrightarrow E$ of the
quotient map. Averaging $s$ over $K$ makes it $K$-equivariant and
leaves it a section. For every such section, the free-extension
module isomorphism is
\[
          A\otimes_\C F\longrightarrow E,\qquad
                               a\otimes f\longmapsto a\,s(f)
\]
\cite[Lemma~2.21]{MCCM}. This is an $A$-module isomorphism;
no tensor-product algebra isomorphism is asserted. The action of
$K$ on $A$ is trivial, so the displayed map is $K$-equivariant.

Taking $H$-invariants yields
$E^H\cong A\otimes F^H$. Averaging over $H$ is exact in
characteristic zero. Applied to the two quotient maps, and averaging
coefficients as above, it identifies $E^H=C$ and $F^H=B$.
Normality of $H$ in $K$ is not needed. Since $z\in R^K$, the
resulting isomorphism is $z$-linear. Quotienting by $z^k$ proves
the Hilbert-series identity.
\end{proof}

\begin{proposition}\label{prop:A-SLP}
The class of $z$ is a strong Lefschetz element of $A$.
\end{proposition}
\begin{proof}
In degree $j$, the dimension of the cokernel of
$z^k:A_{j-k}\to A_j$ is at least $\max(0,h_j-h_{j-k})$.
Equality holds exactly when this multiplication has maximal rank.
Here and below negative graded pieces are zero. Thus
$H_{A/(z^k)}\geq Q_k$ coefficientwise.
For $1\leq k\leq81$, Lemmas~\ref{lem:specialize} and~\ref{lem:free}
and Proposition~\ref{prop:finite} give
\[
 Q_k b\ \leq\ H_{C/(z^k)}\ \leq\ H_{S/(I,z^k)}
                    \ \leq\ H_{S/(J,z^k)}\ =\ Q_k b.
\]
All inequalities are coefficientwise. Hence
$(H_{A/(z^k)}-Q_k)b=0$. The polynomial $b$ has constant
coefficient one and nonnegative coefficients. If the first factor
had a first nonzero coefficient, the product would have that same
nonzero coefficient in its first nonzero degree. Therefore
$H_{A/(z^k)}=Q_k$, proving every required maximal rank for these
$k$. For $k=0$ multiplication is the identity. For $k>81$, every
target with nonzero source is zero because the top degree of $A$
is $81$. Degrees outside the indicated range require no additional
check.
\end{proof}

\begin{proof}[Proof of Theorem~\ref{thm:main}]
For $G_{34}$ the algebra $E$ is a graded Artinian free extension of
$A$ with fiber $F$. Proposition~\ref{prop:A-SLP} gives SLP for $A$,
and Wachi's result for $G_{33}$ gives SLP for $F$
\cite{Wachi25,WachiCode}. Both have symmetric Hilbert functions.
Over a field of characteristic zero, SLP of the base and fiber
therefore implies SLP of the extension
\cite[Proposition~2.22]{MCCM}; see also
\cite[Theorem~6.1]{HW07}. These results allow nonstandard grading,
and their Lefschetz elements lie in degree one. This proves the
$G_{34}$ case. The theorem supplies an element in $E_1$; it does
not assert that $z$ is Lefschetz on $E$ or on $C$.
All other irreducible types are covered by Table~\ref{tab:coverage}.
\end{proof}

\subsection*{Reducible groups and fixed directions}
We now explain why the irreducible result also answers the original
question for every finite complex reflection group.

\begin{corollary}\label{cor:all}
For every finite complex reflection group $W$ on a finite-dimensional
complex vector space, the ordinary coinvariant algebra $E_W$ has the
strong Lefschetz property, including when the representation is
reducible or has a nonzero fixed subspace.
\end{corollary}
\begin{proof}
Average a positive Hermitian inner product over $W$ and decompose
orthogonally into the fixed subspace and nontrivial irreducible summands:
\[
                    V=V^W\oplus V_1\oplus\cdots\oplus V_r.
\]
Each reflection preserves these summands. Since its difference from the
identity has rank one, it acts nontrivially on exactly one summand.
The subgroups generated by reflections supported on the respective
$V_i$ therefore commute, have disjoint supports, and generate $W$.
Thus $W=W_1\times\cdots\times W_r$, with $W_i$ an irreducible
reflection group on $V_i$ and acting trivially on the other summands.

Put $R_i=\Sym(V_i^*)$. Expanding polynomials in the variables of the
other summands shows coefficient by coefficient that invariance under
$W_i$ is precisely invariance of the $R_i$ coefficients. Hence
\[
 R^W=\Sym((V^W)^*)\otimes_\C\bigotimes_{i=1}^r R_i^{W_i}.
\]
Every fixed coordinate is a degree-one invariant and is killed in the
coinvariant quotient. The remaining positive-degree invariant ideal is
generated by the positive-degree invariants in each factor. Consequently
there is a graded algebra isomorphism
\[
 E_W\cong\bigotimes_{i=1}^r
                R_i/((R_i^{W_i})_+R_i)
           =\bigotimes_{i=1}^r E_{W_i}.
\]

Each factor is graded Artinian and has SLP by
Theorem~\ref{thm:main}. The Chevalley--Shephard--Todd polynomial invariant-ring theorem
\cite[Theorem~3]{WachiLecture} and the regular-sequence argument used
above give
\[
 H_{E_{W_i}}(t)=\prod_j(1+t+\cdots+t^{d_{ij}-1}),
\]
where $d_{ij}$ are the basic invariant degrees of $W_i$.
In particular, each factor's Hilbert function is symmetric.
Corollary~2.17 of \cite{MCCM} states that for graded Artinian algebras
with symmetric Hilbert functions in characteristic zero, if $x$ and $y$
are degree-one strong Lefschetz elements of the factors, then
$x\otimes1+1\otimes y$ is strong Lefschetz for their tensor product.
All its hypotheses hold here. Tensor products preserve symmetry of the
Hilbert function, so iteration proves SLP of $E_W$. If $r=0$, the
quotient is $\C$, and its zero degree-one element satisfies the
definition directly.
\end{proof}

\section{Certificates, reproduction, and trusted boundary}\label{sec:trust}
This section describes what the distributed programs establish and
which parts of the theorem remain mathematical inputs or written
arguments. The result is computer-assisted.  There is no complete Lean proof of the theorem.

\subsection*{Certificate format}
The mathematical certificate uses plain-text sparse polynomials in
$a,b,c,d,y,z$ over $\F_{1009}$. The first line is
\texttt{CFG\_POLY\_V1}; the next line gives the prime, variable
count, input count, and number of $G_j$. Records have the forms
\[
\begin{array}{ll}
\texttt{I i polynomial}&\text{the six input polynomials},\\
\texttt{G j polynomial}&\text{a homogeneous nonzero }G_j,\\
\texttt{T i j polynomial}&\text{a nonzero transformation entry}.
\end{array}
\]
Each polynomial is a sum of coefficient--monomial terms with explicit
integer exponents and coefficients reduced modulo $1009$.
Transformation entries are grouped by increasing column, allowing
streamed verification of $I T=G$.
The file \texttt{data/unit-order.txt} specifies the prime, variable names,
positive weights, and weighted degree reverse lexicographic order.
The file \texttt{data/unit-initial.txt} gives candidate leading exponent
tuples; the arithmetic checkers reconstruct them from the $G_j$
instead of accepting them as an independent premise.
The file \texttt{data/unit-starts.txt} gives candidate starts of the
$z$-chains. The monomial checkers reconstruct the complete standard
set and check the cover.

The additional pair certificate starts with
\texttt{CFG\_PAIRS\_V1} and the number of basis polynomials.
It treats every unordered pair exactly once, either by the product
criterion, a chain criterion referring only to earlier accepted
pairs, or a direct S-polynomial reduction. The dependency order
rules out circular chain justifications. A second implementation
audits the entire pair cover and repeats the direct arithmetic.
This verifies Buchberger's criterion, although the proof above needs
only the certified membership and resulting leading-ideal inclusion.

The arithmetic certificate has $534$ membership columns. The pair
cover has $142\,311$ pairs, divided into $1690$ product pairs,
$138\,012$ chain pairs, and $2609$ direct reductions. The two
monomial implementations both cover all $60\,480$ standard
monomials and check all $81$ quotients, including every degree and
the finite tail. The packed-exponent implementation also checks
the degree bounds that prevent exponent overflow.

\subsection*{Reproduction and independence}
The accompanying README gives the requirements and one command that
runs full verification from a clean checkout:
\begin{verbatim}
python3 checks/verify.py --workers 2
\end{verbatim}
The six modular input files are in \texttt{data/inputs/}; the programs
and their entry point are in \texttt{checks/}. The input and basis
records are supplied in a compressed certificate file, together with
the order, leading exponents, and chain starts. The much larger
transformation records are reconstructed locally by the included
Singular recipe before verification. The entry point combines the
supplied input and basis records with these transformation records,
then checks the resulting complete polynomial certificate. No external
certificate download is needed.

The construction uses a Gr\"obner-basis computation to find a witness;
its correctness is not a premise of acceptance. Neither arithmetic
checker imports the producer or performs a basis search. Every
coefficient identity is verified independently, so an incorrect
reconstruction is rejected. The small supplied certificate records
are not overwritten by the construction.

One arithmetic checker uses C++ with FLINT/GMP. The second uses a
separate Python parser and Singular arithmetic, without calling a
basis-search routine. Both independently bind the input records to
the six modular polynomials and verify the coefficientwise
identities. Separate integer implementations check the monomial
cover and quotient Hilbert functions. The exact group model and
the original-coordinate symbolic input comparisons are additional
checks, not consequences of the Gr\"obner computation.
Planted input, basis-coefficient, transformation-coefficient,
leading-term, missing-pure-power, weight, and power-index errors
are rejected by the corresponding checkers. These tests demonstrate
rejection of the specified corruptions; they do not prove universal
checker correctness. Independence here means distinct
implementations, not independent authorship of the programs.

\subsection*{Literature and mathematical trust}
The named group identifications, classification, and invariant
degrees in \eqref{eq:degrees} are imported from
\cite{ST,WachiLecture,OS}. The exact order computation does not
replace those premises. The earlier SLP cases, the tensor-product result, and the two
free-extension results are imported at the locations cited above.
In particular, the $G_{33}$ fiber is not proved anew by these
certificates.

The precise accessible source for Wachi's eight exceptional types
is Wachi's author repository \cite{WachiCode}, whose introduction lists
the proved cases and explicitly states that it has no proof for
$G_{34}$. The abstract of the final article \cite{Wachi25} states
the all-except-one scope. The final article's proof and numbered
theorem were not readable in the source review for this version. In
particular, the final article's treatment of $G_{34}$ remains unchecked;
the 2024 author statement does not establish the scope of every later
version.
No final theorem number is asserted, and the author statement is
not represented as a read proof of the final article. The theorem
here imports that stated result, especially the $G_{33}$ case.
If its scope is corrected, the application above must be reconsidered.

The integral invariant presentation, specialization by relation
matrices, invariant-generation argument, averaging, and
Hilbert-function comparison, together with the product decomposition
and removal of fixed coordinates, are proved in this paper. They are
written mathematical arguments, outside a proof assistant. The
exact execution trusts the parsers, Python integer arithmetic,
SymPy's permutation-group algorithms, FLINT/GMP, Singular, and
the C++ compiler and its runtime. No checker-soundness theorem or
formal proof of the characteristic-zero bridges is claimed.

\subsection*{Further related work and credit}
The invariant construction is that of Conway--Sloane; the subgroup
model and convenient polynomial coordinates follow
Oura--Sekiguchi. Their account also credits Burkhardt's earlier
$G_{33}$ invariants and the work of Orlik, Terao, Enta, Bessis,
and Michel on exceptional invariants. The argument uses their
structural information, with invariant generation also justified
in Proposition~\ref{prop:invariant-generation}.
The transfer method owes its algebraic form to Harima--Watanabe
and McDaniel--Chen--Iarrobino--Macias Marques. Wachi's exceptional
computations and the earlier real, rank-two, and monomial results
are indispensable to the classification-wide conclusion.

Watanabe's survey \cite[\S8]{Watanabe20} discusses the broader
complete-intersection Lefschetz question and distinguishes it from
known reflection cases. Smith's maximal-height work
\cite{Smith24} addresses a related question; the joint McDaniel--Smith
Bott--Samelson work, Smith's announced SLP deduction, and the
McDaniel--Smith--Wachi breakout are discussed in the introduction. Maximal height and
an ambient embedding alone do not imply the all-power ranks proved
here. Kriebel's recent symmetric ideal-chain construction
\cite{Kriebel} concerns a different family and positive-characteristic
behavior; no theorem or code from it is used. Nearby work includes Boij--Migliore--Mir\'o-Roig--Nagel
\cite[Theorem~1.5]{BMMN26} on weak Lefschetz maps in a specified degree
range for equigenerated complete intersections, Nagaoka--Wachi
\cite[Theorem~4]{NW24} on Gorenstein algebras generated by powers of
relative invariants of regular prehomogeneous vector spaces of
commutative parabolic type, and Gui--Sun--Wang--Zhu
\cite[Theorem~1.2]{GSWZ24} on the K\"ahler package for Bott--Samelson
rings of Coxeter systems. Nothing from these three works is used here.
This paper makes no claim of priority over unread later material.

\begin{thebibliography}{99}

\bibitem{BIRS16}
S. Faridi, A. Iarrobino, R.-M. Mir\'o-Roig, L. Smith, and J. Watanabe,
\emph{Lefschetz Properties and Artinian Algebras},
BIRS workshop 16w5114 final report, March 13--18, 2016,
McDaniel's and Smith's talks, pp.~4--6; breakout report, \S4.4, p.~10;
concluding account, \S5, p.~11.
\url{https://archive.birs.ca/files/2016/16w5114/report16w5114.pdf}.

\bibitem{BMMN26}
M. Boij, J. C. Migliore, R. M. Mir\'o-Roig, and U. Nagel,
\emph{Lefschetz Properties of complete intersesctions and maximal variation of hyperplane sections},
October 18, 2026.
\href{https://arxiv.org/abs/2610.25521v1}{arXiv:2610.25521v1}.

\bibitem{GSWZ24}
T. Gui, L. Sun, S. Wang, and H. Zhu,
\emph{On Bott--Samelson rings for Coxeter groups},
first posted August 19, 2024, revised November 5, 2024.
\href{https://arxiv.org/abs/2408.10155v2}{arXiv:2408.10155v2}.

\bibitem{NW24}
T. Nagaoka and A. Wachi,
\emph{The Strong Lefschetz Property of Gorenstein Algebras Generated by Relative Invariants},
March 8, 2024.
\href{https://arxiv.org/abs/2403.05492v1}{arXiv:2403.05492v1}.
\bibitem{CS}
J.~H. Conway and N.~J.~A. Sloane,
\emph{The Coxeter--Todd lattice, the Mitchell group, and related sphere packings},
Math. Proc. Cambridge Philos. Soc. \textbf{93} (1983), 421--440.
\href{https://doi.org/10.1017/S0305004100060746}{doi:10.1017/S0305004100060746}.

\bibitem{HMNW}
T. Harima, J. Migliore, U. Nagel, and J. Watanabe,
\emph{The weak and strong Lefschetz properties for Artinian $K$-algebras},
J. Algebra \textbf{262} (2003), 99--126.
First posted August 26, 2002;
\href{https://arxiv.org/abs/math/0208201v1}{arXiv:math/0208201v1}.

\bibitem{HW06}
T. Harima and J. Watanabe,
\emph{The central simple modules of Artinian Gorenstein algebras},
J. Pure Appl. Algebra \textbf{210} (2007), 447--463.
First posted December 13, 2006;
\href{https://arxiv.org/abs/math/0612344v1}{arXiv:math/0612344v1}.
\href{https://doi.org/10.1016/j.jpaa.2006.10.016}{doi:10.1016/j.jpaa.2006.10.016}.

\bibitem{HW07}
T. Harima and J. Watanabe,
\emph{The strong Lefschetz property for Artinian algebras with non-standard grading},
J. Algebra \textbf{311} (2007), 511--537.
First posted February 15, 2007;
\href{https://arxiv.org/abs/math/0702444}{arXiv:math/0702444}.

\bibitem{Kriebel}
A. Kriebel,
\emph{A symmetric ideal chain with periodic Lefschetz behavior in positive characteristic},
October 22, 2026.
\url{https://zenodo.org/records/23570763}.

\bibitem{McD}
C. McDaniel,
\emph{The strong Lefschetz property for coinvariant rings of finite reflection groups},
J. Algebra \textbf{331} (2011), 68--95.
First posted September 23, 2009;
\href{https://arxiv.org/abs/0909.4184v2}{arXiv:0909.4184v2}.
\href{https://doi.org/10.1016/j.jalgebra.2010.11.007}{doi:10.1016/j.jalgebra.2010.11.007}.

\bibitem{MCCM}
C. McDaniel, S. Chen, A. Iarrobino, and P. Macias Marques,
\emph{Free extensions and Lefschetz properties, with an application to rings of relative coinvariants},
\href{https://arxiv.org/abs/1807.05869v2}{arXiv:1807.05869v2},
first posted July 16, 2018, revised December 7, 2018.
\href{https://doi.org/10.1080/03081087.2019.1598930}{doi:10.1080/03081087.2019.1598930}.

\bibitem{McDSmith}
C. McDaniel and L. Smith,
\emph{Equivariant coinvariant rings, Bott--Samelson rings and Watanabe's bold conjecture},
J. Pure Appl. Algebra \textbf{225} (2021), article 106524.
\href{https://doi.org/10.1016/j.jpaa.2020.106524}{doi:10.1016/j.jpaa.2020.106524}.

\bibitem{OS}
M. Oura and J. Sekiguchi,
\emph{Basic Invariants of the Complex Reflection Group No.34 Constructed by Conway and Sloane},
first posted February 19, 2023.
\href{https://arxiv.org/abs/2302.09695v1}{arXiv:2302.09695v1}.

\bibitem{ST}
G.~C. Shephard and J.~A. Todd,
\emph{Finite unitary reflection groups},
Canad. J. Math. \textbf{6} (1954), 274--304.
\href{https://doi.org/10.4153/CJM-1954-028-3}{doi:10.4153/CJM-1954-028-3}.

\bibitem{Stanley80}
R.~P. Stanley,
\emph{Weyl groups, the hard Lefschetz theorem, and the Sperner property},
SIAM J. Algebraic Discrete Methods \textbf{1} (1980), 168--184.
\href{https://doi.org/10.1137/0601021}{doi:10.1137/0601021}.

\bibitem{Smith24}
L. Smith,
\emph{Notes on Lefschetz Properties and Linear Elements of Maximal Height with Applications},
in \emph{Lefschetz Properties: Current and New Directions},
Springer INdAM Ser. \textbf{59} (2024), 187--210.
\href{https://doi.org/10.1007/978-981-97-3886-1_11}{doi:10.1007/978-981-97-3886-1\_11}.

\bibitem{Stacks}
The Stacks Project Authors, \emph{The Stacks Project},
\href{https://stacks.math.columbia.edu/tag/00N7}{Tag~00N7},
\href{https://stacks.math.columbia.edu/tag/00ND}{Tag~00ND},
\href{https://stacks.math.columbia.edu/tag/02JN}{Tag~02JN}, and
\href{https://stacks.math.columbia.edu/tag/09I3}{Tag~09I3}.

\bibitem{WachiLecture}
A. Wachi, \emph{Ideals generated by invariant polynomials},
in \emph{Proceedings of the 2013 Symposium on Representation Theory},
2013, pp.~140--158, in Japanese, Conjecture~37, Theorem~4, Tables~1--3,
and Lemmas~34--36.
\href{https://doi.org/10.34508/repsympo.2013.0_140}{doi:10.34508/repsympo.2013.0\_140}.
\url{https://www.math.sci.hokudai.ac.jp/coe/sympo/psrt/PSRT_34/14--p140-158.pdf}.

\bibitem{WachiProgram}
\emph{Program of the 2013 Symposium on Representation Theory},
November 26--29, 2013, pp.~159--160; Wachi lecture scheduled
November 29, 2013.
\url{https://www.math.sci.hokudai.ac.jp/coe/sympo/psrt/PSRT_34/15-program-p159-160.pdf}.

\bibitem{WachiAbstract}
A. Wachi, abstract for Functional Analysis talk~12, Mathematical Society of Japan
Spring Meeting, Gakushuin University, March 16, 2014, p.~78;
abstract collection dated March 3, 2014.
\url{https://www.mathsoc.jp/meeting/gakushuin14mar/program/summary14mar_ja_20140303.pdf}.

\bibitem{Wachi25}
A. Wachi,
\emph{The Strong Lefschetz Property for Coinvariant Rings of Finite Complex Reflection Groups},
Experimental Mathematics, first online June 11, 2025.
\href{https://doi.org/10.1080/10586458.2025.2512757}{doi:10.1080/10586458.2025.2512757}.

\bibitem{WachiCode}
A. Wachi, \emph{slp-cpx-refl-gp}, companion programs and author
statement, February 4, 2024, README introduction, revision
\texttt{f2563d1c40ed2efb138bc6e4faaae28a0a19d608}.
\url{https://github.com/wachiakihito/slp-cpx-refl-gp/tree/f2563d1c40ed2efb138bc6e4faaae28a0a19d608}.

\bibitem{Watanabe20}
J. Watanabe, \emph{The Lefschetz properties of graded linear maps}
(in Japanese), Sugaku \textbf{72} (2020), 365--394;
published October 23, 2020.
\href{https://doi.org/10.11429/sugaku.0724365}{doi:10.11429/sugaku.0724365}.
English translation, \emph{The Lefschetz properties},
Sugaku Expositions \textbf{38} (2025), 33--70;
\href{https://doi.org/10.1090/suga/490}{doi:10.1090/suga/490}.

\bibitem{Yoshida}
S. Yoshida,
\emph{The Lefschetz Elements of Coinvariant Algebras of Binary Polyhedral Groups},
Tokyo J. Math. \textbf{36} (2013), 309--336.
\href{https://doi.org/10.3836/tjm/1391177973}{doi:10.3836/tjm/1391177973}.
\end{thebibliography}
\end{document}
